% =============================================================================
%  03_qis.tex —— §3 拟同构
% =============================================================================
\section{拟同构}\label{sec:qi}

\subsection*{拓扑中的拟同构}

§\ref{sec:cx} 已经看到，单纯复形 $X$ 的底拓扑空间 $|X|$（单连通）的一个好不变量，
是单纯链复形 $C\ldot^X$；由 Whitehead 定理，它决定了 $|X|$ 的同伦型。麻烦在于
\emph{函子性}：怎样典范地挑一个三角剖分，从拓扑空间走到复形？标准的数学手腕是
把所有三角剖分放在同等地位上一起考虑，例如让它们彼此都同构。

正如 §\ref{sec:cx} 所说，同一个空间的不同三角剖分（给出不同的链复形 $C\ldot$、
$D\ldot$）之间可能根本没有单纯映射；但取更细的细分、使用单纯逼近，可以找到第三个
链复形 $E\ldot$，嵌入图表
$$
\begin{array}{ccc}
& E\ldot & \\
\swarrow & & \searrow \\
C\ldot & & D\ldot
\end{array}
$$
其中两个映射都是\term{拟同构}{quasi-isomorphism}——它们是同调上诱导同构的链映射。
我们把这个关系扩充成等价关系：两个链复形称为拟同构的，如果它们能由一列上述形状的
拟同构联系起来。于是拟同构的复形有相同的同调，但逆命题\emph{不}成立\yizhuRef{10}；
例如复形
$$
\C[x,y]^{\oplus 2}\xrightarrow{(x,y)}\C[x,y]
\qquad\text{与}\qquad
\C[x,y]\xrightarrow{0}\C
$$
同调相同，却并不拟同构。

拟同构正是我们想在复形上得到的等价关系：由 Whitehead 定理
\eqref{eq:wh}，$|X|$ 与 $|Y|$ 同伦等价，当且仅当存在 $Z$ 诱导拟同构
\begin{equation}\label{eq:dot}
\begin{array}{ccc}
& C^{Z}\ldot & \\
\swarrow & & \searrow \\
C^{X}\ldot & \dashrightarrow & C^{Y}\ldot
\end{array}
\end{equation}
所以我们愿意把拟同构的复形看成同构。它们之间可能并没有映射（正如同一空间的两个
不同三角剖分之间可能没有单纯映射），但我们假装有：手动添上 \eqref{eq:dot} 中
那条虚线箭头，即使它作为真正的映射并不存在。既然底空间终究是一样的，我们就认为
这些复形同构。

特别地，同伦等价是拟同构。代数地看，这就是说：若链映射 $f,g\colon C\ldot\to D\ldot$
之间有一个 $s\colon C\ldot\to D_{\ldot+1}$ 使得 $f-g=\partial_D\circ s+s\circ\partial_C$，
则 $f$、$g$ 在同调上诱导同一个映射。

\begin{yizhu}[“同调相同却不同构”到底难在哪]
若两个复形都同调集中在一个次数上、且由投射模组成，那么它们都拟同构于同一个
“极小”对象；但在一般环上，同调相同并不足以保证拟同构。作者此处所给例子不成立
（见校勘 \ref{kanwu:example}），修正是把系数环换成 $A=\C[\epsilon]/(\epsilon^2)$：
复形 $K=\{A\xrightarrow{\epsilon}A\}$ 的同调是 $H_1=H_0=\C$，零微分复形
$\C\oplus\C[1]$ 的同调也是 $H_1=H_0=\C$，但任何链映射 $K\to\C\oplus\C[1]$ 在
$H_1$ 上都是零（因为 $\Hom_A(\C,A)=0$），反方向也在 $H_1$ 处断掉，故二者不拟同构。
[来源：W. Weibel, \textit{An Introduction to Homological Algebra}, §10.4 例；A 级]
\end{yizhu}

\subsection*{代数几何中的拟同构}

余核映射 \eqref{eq:div}
$$
\{L^{-1}\to\OO_X\}\longrightarrow\OO_D
$$
当然是拟同构（它在层上同调上诱导同构）。同样，任何消解 $F\udot\to\F$
（或 $\F\to F\udot$）给出拟同构 $\{F\udot\}\to\F$（或 $\F\to\{F\udot\}$）。
与拓扑情形一样，我们想把它们看成相等（严格说是同构）$\F\cong\{F\udot\}$，
好让消解的选择变得函子化。于是又需要引入反向箭头
$\{F\udot\}\dashleftarrow\F$（或 $\F\dashleftarrow\{F\udot\}$）。
这正是下面要做的事。

\begin{kanwu}[原文例证不成立：两个复形并不同调]\label{kanwu:example}
原文说 $\C[x,y]^{\oplus2}\xrightarrow{(x,y)}\C[x,y]$ 与
$\C[x,y]\xrightarrow{0}\C$ “同调相同，却不拟同构”。前半句就不对：
前者是 $\C$ 的 Koszul 消解，$H_0=\C$、$H_1=0$；后者的零微分给出
$H_1=\ker(0)=\C[x,y]\ne0$，$H_0=\C$。两者同调并不相同。
真正“同调相同但不拟同构”的标准例子见译注 \ref{yizhu:homology}：
取 $A=\C[\epsilon]/(\epsilon^2)$，比较 $A\xrightarrow{\epsilon}A$ 与
$\C\oplus\C[1]$。原因是 $\Hom_A(\C,A)=0$，$H_1$ 上的同构无法由链映射实现。
建议把原文例子替换为上述例子；本书正文保留原文，另在此处给出判定与修正。
\end{kanwu}

\begin{yizhu}[拟同构不是“同调相同”]\label{yizhu:homology}
$\Hom_A(\C,A)=0$ 的理由：$A$ 中每个元素都被 $\epsilon$ 零化当且仅当它是 $\C$
的倍数；而 $\C$-倍数的元素 $\lambda$ 满足 $\epsilon\lambda=0$，要从 $\C$ 出发的
$A$-模映射必须先穿过 $\C=A/(\epsilon)$，于是像落在 $\epsilon A=0$ 里，只能是零。
因此 $K$ 与 $\C\oplus\C[1]$ 之间没有在 $H_1$ 上为同构的链映射，二者在 $D(A)$
中不同构。这个例子说明：导出范畴里的同构比“同调同构”严格，必须真有链映射（或
屋顶）来实现。
[来源：Weibel 同上；A 级]
\end{yizhu}

\begin{keytakeaways}
\item 拟同构 = 同调上同构的链映射；把全部拟同构可逆化，就得到导出范畴。
\item 拟同构是“正确的同构”：它让消解的选择变得函子化（$\F\cong\{F\udot\}$）。
\item 但“同调相同”不等于“拟同构”——需要真有链映射或屋顶，反例见译注
  \ref{yizhu:homology}。
\end{keytakeaways}
